\chapter{Déploiement}
\label{chap:deployment}

\section{Enjeux du déploiement}
Le déploiement d'une application ne constitue pas une phase purement technique séparée du développement. Il reflète le niveau de maturité du projet, sa reproductibilité et sa capacité à être exécuté dans un environnement stable. Dans le cadre de cette plateforme, le déploiement a été pensé pour répondre à un double besoin :
\begin{itemize}
  \item simplifier la mise en route de l'application dans un contexte de démonstration ou de soutenance ;
  \item poser les bases d'une industrialisation progressive.
\end{itemize}

\section{Conteneurisation avec Docker}
Le projet adopte une approche conteneurisée articulée autour de trois services principaux :
\begin{enumerate}
  \item un conteneur PostgreSQL ;
  \item un conteneur Mailpit pour les tests liés aux e-mails ;
  \item un conteneur applicatif pour l'exécution de Spring Boot.
\end{enumerate}

Cette organisation est particulièrement pertinente pour un projet académique, car elle réduit le coût de configuration locale et garantit une exécution plus homogène entre postes de développement.

\begin{codeblock}[H]
\caption{Fichier \texttt{docker-compose.yml} du projet}
\lstinputlisting[style=academiclst,firstline=1,lastline=39]{\fileDockerCompose}
\end{codeblock}

\section{Docker Compose}
Le fichier \texttt{docker-compose.yml} structure l'environnement en décrivant :
\begin{itemize}
  \item les images ou contextes de build ;
  \item les ports exposés ;
  \item les variables d'environnement ;
  \item les dépendances entre services ;
  \item les volumes persistants.
\end{itemize}

Le recours à Docker Compose améliore fortement la reproductibilité. Il permet de lancer la plateforme avec l'ensemble de ses dépendances sans reconfigurer manuellement la base de données ou les services associés.

\section{Variables d'environnement}
Les variables d'environnement jouent un rôle critique, car elles permettent de séparer configuration et code. Le tableau \ref{tab:variables-env} résume les paramètres essentiels identifiés dans le projet.

\begin{table}[H]
\centering
\caption{Variables d'environnement principales}
\label{tab:variables-env}
\begin{tabularx}{\textwidth}{>{\bfseries}p{3.7cm}p{3cm}X}
\toprule
Variable & Exemple & Rôle \\
\midrule
\texttt{SERVER\_PORT} & \texttt{8080} & port d'exposition de l'application \\
\texttt{DB\_URL} & \texttt{jdbc:postgresql://postgres:5432/parc\_informatique} & URL JDBC de la base \\
\texttt{DB\_USERNAME} & \texttt{postgres} & utilisateur de connexion PostgreSQL \\
\texttt{DB\_PASSWORD} & \texttt{postgres} & mot de passe associé \\
\texttt{JWT\_SECRET} & secret long & signature des jetons JWT \\
\texttt{UPLOAD\_BASE\_PATH} & \texttt{/app/uploads} & emplacement de stockage des fichiers \\
\texttt{APP\_BASE\_URL} & \texttt{http://localhost:8080} & URL racine de l'application \\
\texttt{MAIL\_HOST} & \texttt{mailpit} & service SMTP local \\
\texttt{MAIL\_PORT} & \texttt{1025} & port SMTP \\
\bottomrule
\end{tabularx}
\end{table}

\section{Configuration de production}
La configuration Spring Boot met en évidence plusieurs points intéressants pour un passage futur en production :
\begin{itemize}
  \item la configuration du port et de la base par variables d'environnement ;
  \item la personnalisation du secret JWT ;
  \item la définition des propriétés de mail ;
  \item l'exposition contrôlée d'Actuator et de Swagger ;
  \item l'activation de paramètres JPA et Hikari.
\end{itemize}

\begin{codeblock}[H]
\caption{Configuration applicative Spring Boot}
\lstinputlisting[style=academiclst,firstline=1,lastline=33]{\fileSpringAppYml}
\end{codeblock}

\section{Déploiement applicatif}
Le déploiement local suit une chaîne simple :
\begin{enumerate}
  \item construction ou récupération des images nécessaires ;
  \item initialisation de PostgreSQL ;
  \item démarrage de l'application Spring Boot ;
  \item exposition des services sur les ports définis ;
  \item accès utilisateur viau navigateur.
\end{enumerate}

Cette simplicité est un atout fort dans le cadre d'une soutenance, car elle réduit le risque de défaillance lié à l'environnement.

\section{Perspectives CI/CD}
Le dépôt ne fournit pas, dans son état actuel, une chaîne CI/CD complète industrialisée. Toutefois, la structure du projet permet d'envisager un pipeline simple comprenant :
\begin{itemize}
  \item exécution des tests de contexte ;
  \item vérification du build Maven ;
  \item génération du package \texttt{jar} ;
  \item construction de l'image Docker ;
  \item publication de l'image dans un registre ;
  \item déploiement automatisé vers un environnement cible.
\end{itemize}

Dans une démarche d'amélioration, une intégration GitHub Actions ou GitLab CI serait parfaitement cohérente avec l'architecture retenue.

\section{Sécurité et exploitation en production}
Avant toute mise en production réelle, plusieurs points devraient être renforcés :
\begin{enumerate}
  \item remplacer le secret JWT par une valeur robuste et externalisée ;
  \item durcir la configuration CORS et CSP selon le domaine cible ;
  \item brancher un SMTP réel ;
  \item externaliser durablement le stockage des fichiers ;
  \item ajouter une supervision plus fine des logs et métriques.
\end{enumerate}

\section{Synthèse du chapitre}
Le déploiement de la plateforme a été pensé avec sérieux et pragmatisme. La conteneurisation via Docker, l'usage de Docker Compose et la structuration par variables d'environnement rendent le projet facilement exécutable, démontrable et extensible. La prochaine étape naturelle serait l'automatisation de la chaîne de livraison et le durcissement des paramètres d'exploitation.









